\chapter{Riesgos, calidad y aceptación}

Una oferta profesional no solo enumera beneficios; también muestra los riesgos y cómo se controlan. La matriz siguiente resume los principales riesgos del Sistema AMR-FOD-Kitting y sus mitigaciones.

\begingroup
\small
\setlength{\tabcolsep}{3pt}
\begin{longtable}{@{}P{0.22\textwidth}P{0.08\textwidth}P{0.08\textwidth}P{0.38\textwidth}P{0.12\textwidth}@{}}
\toprule
\textbf{Riesgo} & \textbf{Prob.} & \textbf{Impacto} & \textbf{Mitigación} & \textbf{Resp.}\\
\midrule
\endhead
Interferencia con operarios en pasillos. & Media & Alta & Zonas de seguridad, velocidad limitada, escáner láser, pruebas con peatones y formación. & HSE/Ing.\\
Lecturas RFID/QR incompletas. & Media & Media & Puntos de lectura redundantes, validación HMI y excepciones visibles. & Logística\\
Falsos positivos FOD. & Media & Media & Clasificación asistida, revisión humana y ajuste de iluminación/cámara. & Calidad\\
Mapa desactualizado por cambios de layout. & Media & Alta & Procedimiento de remapeo, control de cambios y versionado de rutas. & Ing.\\
Saturación de red WiFi. & Baja & Alta & Estudio de cobertura, QoS industrial y modo degradado con cola local. & IT/OT\\
Baja aceptación del operario. & Media & Alta & Piloto participativo, HMI simple, formación corta y medición de carga operativa. & Producción\\
Coste percibido superior al beneficio inicial. & Media & Media & Piloto medido, KPIs antes--después y expansión condicionada a evidencia. & Dirección\\
\bottomrule
\end{longtable}
\endgroup

\section{Criterios de aceptación por fase}

\begin{itemize}
  \item \textbf{Seguridad:} cero contactos no previstos, parada de emergencia verificada y rutas validadas con operarios.
  \item \textbf{Operación:} al menos 95\% de misiones completadas en piloto, con causas de fallo clasificadas.
  \item \textbf{Trazabilidad:} al menos 98\% de misiones con registro completo de identificador, destino, hora y confirmación.
  \item \textbf{FOD:} rondas registradas y eventos revisables por calidad, sin sustituir inspecciones obligatorias.
  \item \textbf{Escalabilidad:} gestor de flota probado antes de pasar de 1 a 7 robots.
\end{itemize}

\section{Gobierno de datos}

Los datos de misiones deben conservarse con identificador único, sello temporal, estación, robot, operador de confirmación, tipo de misión, excepciones y evidencias FOD. Se recomienda un periodo mínimo de conservación de 12 meses para análisis de mejora continua y auditorías internas. El acceso debe separarse por perfiles: operario, supervisor, mantenimiento, calidad y administrador.
